\documentclass{statsoc}

\usepackage[a4paper]{geometry}
\usepackage{graphicx}
\usepackage{amsmath,amssymb}
%\usepackage{longtable,booktabs}
\usepackage{natbib}
\usepackage{hyperref}

\usepackage[utf8]{inputenc}

\title[Ensemble Prediction for CD]{Ensemble Prediction of Time to Event Outcomes with Competing Risks: A Case Study of Surgical Complications in Crohn's Disease: SUPPLEMENTARY MATERIAL}
\author{Michael C Sachs}
\address{Department of Medicine, Karolinska Institutet
Stockholm,
Sweden.}
\email{michael.sachs@ki.se}
\author{Andrea Discacciati}
\address{Institute of Environmental Medicine, Karolinska Institutet,
Stockholm, Sweden.}
\author{Åsa H Everhov}
\address{Department of Medicine, Karolinska Institutet,
Stockholm, Sweden.}
\author{Ola Olén}
\address{Department of Medicine, Karolinska Institutet,
Stockholm, Sweden.}
\author[Michael C Sachs {\it et al.}]{Erin E Gabriel}
\address{Department of Medical Epidemiology and Biostatistics, Karolinska Institutet,
Stockholm, Sweden.}


\begin{document}
  

\section{Detailed specification of simulation study}

\subsection{Data generation mechanisms}

Let $N(\mu, \sigma)$ denote the normal distribution with mean $\mu$ and standard deviation $\sigma$. For each replicate of the simulation study, and for each of $i = 1, \ldots, 500$ independent subjects, a vector $\mathbf{X_i} = (X_{i1}, \ldots, X_{i20})$ of predictors was generated as follows: 

\begin{enumerate}
\item $X_{ij} \sim N(0, 0.1)$ for $j = 1, \ldots, 5$
\item $X_{ij} = \sum_{k = 1}^5 0.25 * X_{ik} + \varepsilon_{i}(0.1)$, for $j = 6, \ldots, 10$, where $\varepsilon_{i}(0.1) \sim N(0, 0.1)$
\item $X_{ij} = \sum_{k = 6}^{10} 0.15 * X_{ik} + \varepsilon_{i}(0.5)$, for $j = 11, \ldots, 15$, where $\varepsilon_{i}(0.5) \sim N(0, 0.5)$
\item $X_{ij} = \sum_{k = 11}^{16} 0.05 * X_{ik} + \varepsilon_{i}(0.65)$, for $j = 16, \ldots, 20$, where $\varepsilon_{i}(0.65) \sim N(0, 0.65)$
\end{enumerate}

This induces a correlation structure among the predictors. Then for each scenario the models for the scale parameter of a Weibull distribution are as follows: 

\begin{itemize}
    \item[Scenario 0:] $\gamma_{i1} = \exp{(\varepsilon_{i}(0.35))}$ where $\varepsilon_{i}(0.35) \sim N(0, 0.35)$.
    \item[Scenario A:] $\gamma_{i1} = \exp{(-2 + 2.5 * X_{i1})}$.
    \item[Scenario B:] $\gamma_{i1} = \exp{(6 + \mathbf{Z_{ib}}^T \beta_b)}$ where $\mathbf{Z_{ib}} = (X_{i1}, X_{i6}, X_{i1} * X_{i6}, b_3[X_{i1}], b_3[X_{i6}])$, and where $b_3$ denotes a b-spline basis expansion of degree 3 with 3 knots, and $\beta_b = (.75, .5, 6.1, 1.02, -2.03, 1, 2, 1, .6, .1, 3)$.
    \item[Scenario C:] $\gamma_{i1} = \exp{(\mathbf{Z_{ic}}^T \beta_c)}$, where $\mathbf{Z_{ic}} = (X_{i1}, X_{i6}, X_{i11}, X_{i16}, X_{i20}, X_{i1} * X_{i6}, \cos[X_{i11} / .1], X_{i16} * I[X_{i16} < 0])$, where $I(\cdot)$ denotes the indicator function that equals 1 if the condition is true, and 0 otherwise, and where $\beta_c = (1.1, 1.4, -2.1, -1.2, -2.3, -1.5, 6.7, .5) / 4$. 
    \item[Scenario D:] $\gamma_{i1}$ generated as in Scenario C, and additionally $\delta_{i1} = \exp((X_{i1}, X_{i6}, X_{i11}, X_{i16},X_{i20})^T \alpha_d)$ where $\alpha_d = (1.1, 1.4, -2.1, -1.2, -2.3, -1.5, 6.7, .5) / 3$.
    \item[Scenario E:] The real data example was used and bootstrapping with replacement was done to sample datasets for the simulation. 
\end{itemize}

Then, with $t_{\star} = 26.5$, $\kappa_1 = \exp(2.5 * X_{i2})$, $\kappa_2 = 2.5$, and $\gamma_2 = 3.5$, we create the rescaling terms: 

$$C_{i\gamma} = \frac{t_{\star}}{(-\log(0.8))^{-\gamma_2}}$$
and
$$C_{i\kappa} = \frac{t_{\star}}{(-\log(0.93))^{-\kappa_2}}$$ 

which are then applied as follows: 

$$
\gamma^*_{i1} = \gamma_{i1} * n^{-1}\sum \frac{C_{i\gamma}}{\gamma_{i1}}
$$
and 
$$
\kappa^*_{i1} = \kappa_{i1} * n^{-1}\sum \frac{C_{i\kappa}}{\kappa_{i1}}. 
$$

Then $Y_{i1}$ (time to surgery) is sampled from a Weibull distribution with scale $\gamma_{i1}^*$ and shape $\gamma_2$, and $Y_{i2}$ (time to death) is sampled from a Weibull distribution with scale $\kappa_{i1}^*$ and shape $\kappa_2$. The rescaling ensures that the average cumulative incidence at $t_{\star}$ is 0.2 for surgery, and 0.07 for death. 

For scenario D, censoring times $Y_{i0}$ were generated by a Weibull distribution with scale $\delta_{i1}$ and shape 1, and otherwise censoring times were uniformly distributed between the 5\% percentile of the distribution of event times and its maximum. 

For each subject, the true cumulative incidence at $t_{\star}$ for cause 1 was calculated as

\begin{equation*}
\begin{split}
P(Y_{i1} < t_{\star} \mbox{ and } Y_{i1} < Y_{i2}) = E\{P(Y_{i1} < t_{\star} \mbox{ and } Y_{i1} < x | Y_{i2} = x)\} \\
= \int_{t_\star}^\infty P(Y_{i1} < t_{\star} \mbox{ and } Y_{i1} < x) dF(x; \kappa_{i1}, \kappa_{i2}) \\
 = \int_{0}^{t_{\star}} F(x; \gamma_{i1}, \gamma_{i2}) dF(x; \kappa_{i1}, \kappa_{i2})  + 
\int_{t_{\star}}^\infty F(t_{\star}; \gamma_{i1}, \gamma_{i2}) dF(x; \kappa_{i1}, \kappa_{i2}) 
\end{split}
\end{equation*}
where $F(x; a_1, a_2)$ is the cumulative distribution function of the Weibull distribution with scale $a_1$ and shape $a_2$. 


\subsection{Analysis}

For each simulated dataset, we calculated pseudo-observations at a fixed  grid of time points $\{17.5, 20, 26.5, 35\}$, where the time point of interest is 26.5, using the \texttt{pseudo} package. A validation dataset was generated independently with pseudo observations calculated only at the time point of interest. 

The SuperLearner model was run for the pseudo-observations using the entire grid of time points, with time as an additional covariate and again with only the time point of interest. The following learning algorithms from the \texttt{SuperLearner} R package were used in the library with their default settings: 

\begin{verbatim}
> SL.library
[[1]]
[1] "SL.glm"      "screen.corP"
[[2]]
[1] "SL.step"
[[3]]
[1] "SL.gam"      "screen.corP"
[[4]]
[1] "SL.ksvm"
[[5]]
[1] "SL.ranger"
[[6]]
[1] "SL.rpart"
[[7]]
[1] "SL.glmnet"
[[8]]
[1] "SL.polymars" "screen.corP"
[[9]]
[1] "xgb_200_2_0.01"
[[10]]
[1] "xgb_200_2_0.1"
[[11]]
[1] "xgb_200_2_0.2"
\end{verbatim}

The combination of learners was optimized to maximize the time varying AUC, as described in the manuscript, with a fixed $\lambda = 100$. 

The SuperLearner for the binary outcome was run using inverse probability of censoring weights estimated by the Kaplan-Meier estimate of not being censored at the minimum of the observed time and 26.5 (the time of interest). The learners in that library were the following: 

\begin{verbatim}
[[1]]
[1] "SL.glm"
[[2]]
[1] "SL.gam"
[[3]]
[1] "SL.ksvm"
[[4]]
[1] "SL.ranger"
[[5]]
[1] "SL.rpart"
[[6]]
[1] "SL.glmnet"
[[7]]
[1] "SL.polymars"
[[8]]
[1] "xgb_200_2_0.01"
[[9]]
[1] "xgb_200_2_0.1"
[[10]]
[1] "xgb_200_2_0.2"
\end{verbatim}

The weights were passed to the SuperLearner function, which passes them to the methods individually, when they are supported by that method. The combination of learners was optimized to minimize the weighted non-negative log likelihood, via \texttt{method.NNloglik}. 
The CoxBoost and survival random forests predictive models were estimated using the \texttt{CoxBoost::CoxBoost} and \texttt{randomForestsSRC::rfsrc} functions with their default settings, respectively. Predictions of the cumulative incidence function at the time point of interest were obtained by the predict methods from those packages. 

The simulations were run in R version 3.3.1 on a Windows 2008 server. 

R code for data analysis and simulations is available as an R source package at the following link. The code corresponding to what is reported in the manuscript is tagged with release date 2018-12-24. 

\href{http://github.com/sachsmc/pseupersims}{http://github.com/sachsmc/pseupersims}

\section{Description of Crohn's Disease Study Population}

\begin{table}
\caption{\label{tab1} Description of study population and potential predictors. Demographics and family characteristics. } 
\centering
\begin{tabular}{p{3cm}rrrrr}
\hline
& Censored & Surgery & Died & Emigrated &  Overall \\
& (n=7548) & (n=1397) & (n=567) & (n=93) & (n=9605) \\
\hline
\textbf{Age at onset} & & & & & \\
Mean (SD) & 40.9 (16.8) & 43.2 (18.0) & 70.9 (14.1) & 33.4 (12.0) & 42.9
(18.2) \\
Median {[}Min, Max{]} & 38 {[}18, 93{]} & 41 {[}18, 91{]} &
74 {[}19, 98{]} & 32 {[}18, 79{]} & 40 {[}18, 98{]} \\
\textbf{Male} & & & & & \\
Yes & 3553 (47.1\%) & 715 (51.2\%) & 258 (45.5\%) & 43 (46.2\%) & 4569
(47.6\%) \\
No & 3995 (52.9\%) & 682 (48.8\%) & 309 (54.5\%) & 50 (53.8\%) & 5036
(52.4\%) \\
\textbf{Inpatient} & & & & & \\
Yes & 1856 (24.6\%) & 540 (38.7\%) & 305 (53.8\%) & 24 (25.8\%) & 2725
(28.4\%) \\
No & 5692 (75.4\%) & 857 (61.3\%) & 262 (46.2\%) & 69 (74.2\%) & 6880
(71.6\%) \\
\multicolumn{3}{l}{\textbf{Highest parental education}} & & & \\
Primary school & 1211 (16.0\%) & 260 (18.6\%) & 82 (14.5\%) & 0 (0\%) &
1553 (16.2\%) \\
Secondary school & 356 (4.7\%) & 76 (5.4\%) & 7 (1.2\%) & 4 (4.3\%) &
443 (4.6\%) \\
Upper secondary school 2 years & 1888 (25.0\%) & 313 (22.4\%) & 38
(6.7\%) & 6 (6.5\%) & 2245 (23.4\%) \\
Upper secondary school 3 years & 2387 (31.6\%) & 473 (33.9\%) & 428
(75.5\%) & 64 (68.8\%) & 3352 (34.9\%) \\
University \textless{} 3 years & 733 (9.7\%) & 119 (8.5\%) & 7 (1.2\%) &
10 (10.8\%) & 869 (9.0\%) \\
University 3 years or more & 901 (11.9\%) & 148 (10.6\%) & 5 (0.9\%) & 8
(8.6\%) & 1062 (11.1\%) \\
Post graduate & 72 (1.0\%) & 8 (0.6\%) & 0 (0\%) & 1 (1.1\%) & 81
(0.8\%) \\
\multicolumn{3}{l}{\textbf{Mother born outside Sweden}} & & & \\
Yes & 919 (12.2\%) & 150 (10.7\%) & 30 (5.3\%) & 13 (14.0\%) & 1112
(11.6\%) \\
No & 6629 (87.8\%) & 1247 (89.3\%) & 537 (94.7\%) & 80 (86.0\%) & 8493
(88.4\%) \\
\multicolumn{3}{l}{\textbf{Born outside Sweden}} & & & \\
Yes & 1184 (15.7\%) & 220 (15.7\%) & 83 (14.6\%) & 62 (66.7\%) & 1549
(16.1\%) \\
No & 6364 (84.3\%) & 1177 (84.3\%) & 484 (85.4\%) & 31 (33.3\%) & 8056
(83.9\%) \\
\hline
\end{tabular}
\end{table}


\begin{table}
\caption{Description of study population and potential predictors. Medications and IBD characteristics. \label{tab1b}} 
\centering
\begin{tabular}{lrrrrr}
\hline
& Censored & Surgery & Died & Emigrated &  Overall \\
& (n=7548) & (n=1397) & (n=567) & (n=93) & (n=9605) \\
\hline
\textbf{Immunomodulators} &  1202 (15.9\%) & 143 (10.2\%) & 39 (6.9\%) & 9 (9.7\%) & 1393 (14.5\%) \\
\textbf{anti-TNF}  & 115 (1.5\%) & 13 (0.9\%) & 0 (0\%) & 0 (0\%) & 128 (1.3\%) \\
\textbf{Infliximab}  & 44 (0.6\%) & 4 (0.3\%) & 0 (0\%) & 0 (0\%) & 48 (0.5\%) \\
\textbf{Cortico steroids} & 3301 (43.7\%) & 524 (37.5\%) & 192 (33.9\%) & 20 (21.5\%) & 4037 (42.0\%) \\
\textbf{Local steroids} & 1824 (24.2\%) & 309 (22.1\%) & 77 (13.6\%) & 13 (14.0\%) & 2223 (23.1\%) \\
\textbf{Rectal ASA}  & 2181 (28.9\%) & 239 (17.1\%) & 96 (16.9\%) & 22 (23.7\%) & 2538 (26.4\%) \\
\textbf{Systemic ASA}  & 2579 (34.2\%) & 286 (20.5\%) & 113 (19.9\%) & 23 (24.7\%) & 3001 (31.2\%) \\
\textbf{Antibiotics}  & 2035 (27.0\%) & 408 (29.2\%) & 150 (26.5\%) & 11 (11.8\%) & 2604 (27.1\%) \\
\textbf{Extraintestinal Manifestations} & 172 (2.3\%) & 21 (1.5\%) & 8 (1.4\%) & 1 (1.1\%) & 202 (2.1\%) \\
\textbf{Primary Sclerosing Cholangitis}  & 29 (0.4\%) & 1 (0.1\%) & 6 (1.1\%) & 1 (1.1\%) & 37 (0.4\%) \\
\textbf{Relative has UC}  & 503 (6.7\%) & 89 (6.4\%) & 18 (3.2\%) & 3 (3.2\%) & 613 (6.4\%) \\
\textbf{Relative has CD}  & 591 (7.8\%) & 117 (8.4\%) & 17 (3.0\%) & 4 (4.3\%) & 729 (7.6\%) \\
\textbf{Stricturing/penetrating} & 6979 (92.5\%) & 1263 (90.4\%) & 515 (90.8\%) & 85 (91.4\%) & 8842 (92.1\%) \\
\textbf{Ileocecal location}  & 3236 (42.9\%) & 616 (44.1\%) & 289 (51.0\%) & 46 (49.5\%) & 4187 (43.6\%) \\
\textbf{Perianal disease}  & 457 (6.1\%) & 104 (7.4\%) & 10 (1.8\%) & 3 (3.2\%) & 574 (6.0\%) \\
\hline
\end{tabular}
\end{table}


\begin{table}
\caption{Description of study population and potential predictors. Comorbidities. \label{tab1c}} 
\centering
\begin{tabular}{lrrrrr}
\hline
& Censored & Surgery & Died & Emigrated &  Overall \\
& (n=7548) & (n=1397) & (n=567) & (n=93) & (n=9605) \\
\hline
\textbf{Diabetes} & 297 (3.9\%) & 46 (3.3\%) & 95 (16.8\%) & 1 (1.1\%) & 439 (4.6\%) \\
\textbf{Hypertension} & 642 (8.5\%) & 120 (8.6\%) & 211 (37.2\%) & 2 (2.2\%) & 975 (10.2\%) \\
\textbf{Ischemic heart disease} & 264 (3.5\%) & 50 (3.6\%) & 154 (27.2\%) & 1 (1.1\%) & 469 (4.9\%) \\
\textbf{Cerebrovascular disease} & 131 (1.7\%) & 35 (2.5\%) & 81 (14.3\%) & 0 (0\%) & 247 (2.6\%) \\
\textbf{Congestive heart disease} & 95 (1.3\%) & 23 (1.6\%) & 138 (24.3\%) & 0 (0\%) & 256 (2.7\%) \\
\textbf{COPD}  & 118 (1.6\%) & 27 (1.9\%) & 85 (15.0\%) & 0 (0\%) & 230 (2.4\%) \\
\textbf{Kidney failure} & 107 (1.4\%) & 12 (0.9\%) & 54 (9.5\%) & 0 (0\%) & 173 (1.8\%) \\
\hline
\end{tabular}
\end{table}

%\clearpage


%\section{Bibliography}

%\bibliographystyle{rss}
%\bibliography{irisbiblio}

\end{document}